\documentclass{article}
\usepackage[utf8]{inputenc}
\usepackage[T1]{fontenc}
\usepackage{geometry}
\geometry{
    a4paper,
    left=2.5cm,
    right=2.5cm,
    top=2.5cm,
    bottom=2.5cm
}
\usepackage{graphicx}
\usepackage{float}
\usepackage{hyperref}
\usepackage[english]{babel}
\usepackage{amsmath}
\usepackage{amssymb}
\usepackage{array}
\usepackage{booktabs}
\usepackage{caption}
\usepackage{subcaption}
\usepackage{enumitem}
\usepackage{fancyhdr}
\usepackage{titlesec}
\usepackage{url}
\usepackage{textcomp}
\usepackage{xcolor}
\usepackage{listings}
\usepackage{multirow}
\usepackage{longtable}
\usepackage{tabularx}
\usepackage{times}

% Set up page style
\pagestyle{fancy}
\fancyhf{}
\rfoot{\thepage}
\renewcommand{\headrulewidth}{0pt}

% Title formatting
\titleformat{\chapter}[display]{\normalfont\bfseries}{\chaptertitle\thechapter}{0pt}{\Huge}
\titlespacing*{\chapter}{0pt}{50pt}{40pt}

% Document info
\title{Template Document}
\author{Academic Support Agent}
\date{\today}

\begin{document}
\maketitle

\begin{document} % All begin commands must be paired with an end command somewhere

	\maketitle % creates title using information in preamble (title, author, date)
    
	%New section is created
	\section{Introduction}
	%Plain text is just written directly in this document like this:
	This is a template. You can make references to things like \citep{latexsource}, which will show up in the bibliography (\sectref{bibsect}.)

	%New section about equations
	\section{Equations}
	Here is an inline equation: \ieq{F=ma}. Below is a block equation which we can label as \eqref{example_eqn}.

	\begin{equation} 
		\label{example_eqn}
		a^2+b^2=c^2
	\end{equation}

	%New section about tables and figures
	\section{Figures and tables}
	\figref{figtest} is an example figure. 

	\begin{figure}[h!]
  		\includegraphics[width=\linewidth]{images/example.png}
  		\caption{An example image.}
  		\label{figtest}
	\end{figure}


	\tableref{tabletest} is an example table.

	\begin{table}[h!]
		\caption{An example table.}
		\label{tabletest}
		\begin{center}
  			\begin{tabularx}{\textwidth}{ X X X }
    				\toprule
    				Heading 1 & Heading 2 & Heading 3 \\ 
    				\midrule
    				1 & 2 & 3 \\ 
    				4 & 5 & 6 \\ 
    				7 & 8 & 9 \\
    				\bottomrule
  			\end{tabularx}
		\end{center}
	\end{table}

	%New section about code
	\section{Code}
	Below is some C++ code.
	\begin{lstlisting}[language=C++]
tmp<fvVectorMatrix> tUEqn
(
	fvm::div(phi, U)
	+ MRF.DDt(U)
	+ turbulence->divDevReff(U)
	==
	fvOptions(U)
);
\end{lstlisting}
   
	
	%Add bibliography at the end of document	
	
	\label{bibsect}
\end{document}